\documentclass[11pt]{article}
\usepackage[margin=1in]{geometry}
\usepackage{amsmath,amssymb,amsthm}
\usepackage[colorlinks=true,linkcolor=blue,citecolor=blue,urlcolor=blue]{hyperref}

\newcommand{\R}{\mathbb R}
\newcommand{\Q}{\mathbb Q}
\newcommand{\PF}{\operatorname{PF}}
\newcommand{\F}{\mathcal F}

\title{Literature-Implied Partial Answer: Real Parseval Frame Connectedness}
\author{Run fa\_banach\_001, lane 8}
\date{2026-07-18}

\begin{document}
\maketitle

\paragraph{Status.}
This packet records a \emph{literature-implied partial subcase} of the real frame connectedness problem extracted from Needham--Shonkwiler, \emph{Admissibility and Frame Homotopy for Quaternionic Frames}, arXiv:2108.02275. It is not claimed as a new proof.

\paragraph{Source problem.}
The source paper asks whether one can interpolate between arbitrary frames with fixed frame spectrum and fixed vector norms. It gives complete quaternionic answers, recalls that the complex version is known, and states in the introduction that for real frames with nonconstant spectrum or norms the connectedness question remains open. Its discussion formulates the remaining real issue as the challenge of characterizing which data \((\lambda,r)\) make
\[
\F_{\lambda}^{\R^d,N}(r)
\]
connected.

\paragraph{Supporting theorem.}
Caine--Needham--Shonkwiler, \emph{Optimization and the Topology of Spaces of Parseval Frames}, arXiv:2505.14860, studies the Parseval prescribed-norm subcase
\[
\PF_d^\R(r)=\{F\in \R^{d\times n}: FF^*=I_d,\ \|f_i\|^2=r_i\}.
\]
For an admissible \(r\in \Q_+^n\), write \(r_{(1)}\ge \cdots \ge r_{(n)}\) and
\[
k_\ell(r)=\min\left\{k:\sum_{i=1}^k r_{(i)}>\ell\right\},
\qquad 1\le \ell\le d-1.
\]
Theorem 4.5 of arXiv:2505.14860 says, among other things, that if there is a constant \(c\) with
\[
k_\ell(r)\ge c\,\frac{n\ell}{d}\quad (1\le \ell\le d-1)
\]
and, for real frames,
\[
c\ge \frac{d}{n}\left(\frac{q+d+1}{d-1}\right),
\qquad d\ge q+2,
\]
then \(\PF_d^\R(r)\) is \(q\)-connected. Taking \(q=0\) gives a path-connectedness criterion.

\paragraph{Concrete connectedness family.}
Corollary 4.9 of arXiv:2505.14860 gives the following open-family version. If \(d\ge 2\), \(r\in \Q_+^n\) is admissible, and the above constant \(c\) satisfies
\[
c\ge \frac{d}{n}\left(\frac{2d}{d-1}\right),
\]
then there exists \(\epsilon>0\) such that every admissible rational vector
\[
s\in \Q_+^n,\qquad |s_i-r_i|<\epsilon\quad(1\le i\le n),
\]
has \(\PF_d^\R(s)\) path-connected. In particular, for the equal-norm vector \(r_i=d/n\), where one may take \(c=1\), this gives an open rational neighborhood of connected real Parseval frame spaces whenever
\[
n\ge \frac{2d^2}{d-1}.
\]

\paragraph{Why this answers a subcase.}
The source problem includes fixed-spectrum real frame spaces. The Parseval condition is the fixed-frame-operator subcase with spectrum \((1,\ldots,1)\), and prescribed squared norms \(r\). Thus Theorem 4.5 and Corollary 4.9 give substantial sufficient conditions under which a nonconstant real prescribed-norm subfamily is path-connected. The implication is agent-identified from the two papers; arXiv:2505.14860 cites arXiv:2108.02275 as part of the frame-homotopy literature, but does not present itself as a complete answer to the 2108.02275 residual real characterization problem.

\paragraph{Limitations and remaining open part.}
This does not settle the full problem. The result is restricted to \(S=I_d\), rational prescribed norms, and sufficient connectedness conditions. The supporting paper explicitly notes that a complete characterization of the vectors \(r\) for which \(\PF_d^\R(r)\) is path-connected remains open. It also separately records known planar polygon quotient disconnection criteria, while warning that a quotient by the disconnected sign-change group does not by itself give the converse for the frame space.

\paragraph{Evidence checked.}
The run indexes were searched for arXiv ids 2108.02275 and 2505.14860 and for frame-space, Parseval-frame, polygon, and real connectedness keywords. Before this packet, the only relevant durable packet was the lane-8 original partial result reducing \(d=2\) real fixed-spectrum frame spaces to a labelled polygon-level double cover. The 2505.14860 paper is a separate later literature source and is therefore recorded here under \texttt{solutions/literature\_implied\_answers/}.

\paragraph{References.}
\begin{itemize}
\item T. Needham and C. Shonkwiler, \emph{Admissibility and Frame Homotopy for Quaternionic Frames}, arXiv:2108.02275; Linear Algebra Appl. 645 (2022), 237--255.
\item A. Caine, T. Needham, and C. Shonkwiler, \emph{Optimization and the Topology of Spaces of Parseval Frames}, arXiv:2505.14860, 2025.
\end{itemize}

\end{document}
